\section{The half-residue strict row for all three pure cubic labels}
\label{sec:pure-cubic-half-residue-strict}

\paragraph{Inputs and scope.}
The only non-elementary class input in this proof is the already
proved conclusion of AC-G152-class for its degree-twelve field.
An explicit degree-two subfield and two units from the same exact
data suffice.  No new analytic class certificate or complete unit
basis is asserted.  The trivial cubic character is included;
its extra class and unit arguments are supplied below.

Put $p=269$, $K=\mathbb Q(\omega)$, $t=\sqrt[3]p>0$, and let
$\chi:G_K\to\mu_3$ describe $K(t)/K$.
Write $c_p$ for the actual cyclotomic character and $\omega_p$
for its Teichmuller part.  For a discrete module $D$ set
\[
 \begin{gathered}
 S_{\rm str}(K,D)=\ker\bigl(H^1_{\rm cts}(G_K,D)
                                  \longrightarrow\mathcal L(D)\bigr),\\
 \mathcal L(D)=\prod_{v\mid p}H^1(G_{K_v},D)
           \times\prod_{v\nmid p,\infty}H^1(I_v,D).
 \end{gathered}
 \tag{PCH.1}
\]

\begin{theorem}[AC-G154: the pure cubic half-residue strict rows]
\label{AC-G154}
Let $E/\mathbb Q_p$ be finite and let $O$ be its original ring of
integers, containing the values of $\chi$.
Let $j\in\{0,1,2\}$, $m\in\mathbb Z$, $e\geq1$, and let
$\kappa:G_K\to1+\mathfrak m_O$ be continuous and unramified at
every finite place away from $p$.
If $e+m\equiv133\pmod{268}$, then
\[
 \begin{gathered}
 D=(E/O)(\chi^{-j}\omega_p^{-m}\kappa^{-1}c_p^{-e}),\\
                         S_{\rm str}(K,D)=0.
 \end{gathered}
 \tag{PCH.2}
\]
The corresponding residual strict groups vanish over every finite
field containing the character values.
No restriction on the absolute ramification index of $O$ is imposed.
For $m=0,\kappa=1$ the finite labels are exactly
$1,\chi,\chi^2$, already present on the actual AC-K25 hosts.
\end{theorem}

\begin{proof}
\proofstep{1. An inherited degree-six field and the residual quotient.}
Use the real quartic field and degree-twelve extension of
AC-G152-class, choosing the sign of its quadratic generator so that
\[
 \begin{gathered}
 F=\mathbb Q(w),\qquad w^4-pw^2+25p=0,\\
 w^2=(p+13s)/2,\qquad s^2=p,\\
 L^*=\mathbb Q(t,w),\qquad L_0=\mathbb Q(t,s).
 \end{gathered}
 \tag{PCH.3}
\]
Set $z=s/t$.  The exact identities
\[
       s=(2w^2-p)/13,\qquad z^2=t,\qquad z^3=s,
       \qquad z^6=p
 \tag{PCH.4}
\]
give $L_0=\mathbb Q(z)$.
The polynomial $X^6-p$ is Eisenstein at $p$, so $[L_0:\mathbb Q]=6$.
The retained class theorem gives
$[L^*:\mathbb Q]=12$ and $\operatorname{Cl}(L^*)[p]=0$.
Consequently $[L^*:L_0]=2$.
Class extension followed by norm is multiplication by two, hence
\[
                         \operatorname{Cl}(L_0)[p]=0.
 \tag{PCH.5}
\]
This is an inherited class-primality statement, with no new
class-number computation.

Put $M_0=KL_0=K(t,s)$ and $N=K(s)$.
The normal closure $K(t)$ has group $S_3$ and unique quadratic
subfield $K$.  Since $\mathbb Q(s)$ is real, their intersection
is $\mathbb Q$.  Thus
\[
 \begin{gathered}
 \operatorname{Gal}(M_0/\mathbb Q)=S_3\times C_2,
 \qquad [M_0:\mathbb Q]=12,\\
 \Delta=\operatorname{Gal}(M_0/K)=C_3\times C_2.
 \end{gathered}
 \tag{PCH.6}
\]
Let $c(t)=\omega t$, $c(s)=s$, and let $b(s)=-s$ with $b(t)=t$.
They fix $K$, commute, and satisfy
\[
                         c(z)=\omega^{-1}z,\qquad b(z)=-z.
 \tag{PCH.7}
\]

Write $k=O/\mathfrak m_O$ and let $a$ be the residual character
in (PCH.2).  Since $\kappa$ reduces to one, its Kummer quotient is
\[
 \begin{aligned}
 \eta=a/\bar c_p
     =\bar\chi^{-j}\bar c_p^{-(e+m+1)}
     &=\bar\chi^{-j}\bar\tau_p,\\
 \tau_p&=\omega_p^{134}.
 \end{aligned}
 \tag{PCH.8}
\]
The quadratic character $\tau_p$ cuts out $\mathbb Q(\sqrt p)$,
because $p\equiv1\pmod4$.
Thus $M_0$ splits $\eta$ for all three values of $j$.
For $j=0$ its kernel contains $C_3$; the descent below explicitly
allows this nonfaithful prime-to-$p$ cover.
In every case $\eta^{-1}(b)=-1$.
The factor $\tau_p$ belongs only to this residual splitting field,
not to the original finite host label.

\proofstep{2. The extra class argument for the trivial cubic label.}
We first prove $\operatorname{Cl}(N)[p]=0$ using its three
quadratic subfields.  For $K$, the imaginary quadratic Minkowski
bound $2\sqrt3/\pi<2$ gives an integral ideal representative of
norm one, so its class group is trivial.
For $\mathbb Q(s)$, extension to $L^*$ followed by norm is
multiplication by six.  Since $p\nmid6$, (PCH.3) and
AC-G152-class give $\operatorname{Cl}(\mathbb Q(s))[p]=0$.

The third subfield is $\mathbb Q(\sqrt{-807})$.
Its fundamental discriminant is $-807$, and its Minkowski norm
bound satisfies
\[
        \frac{2\sqrt{807}}\pi<\frac{2\sqrt{807}}3<19,
        \qquad 4\cdot807=3228<3249=9\cdot19^2.
 \tag{PCH.9}
\]
In any quadratic field the number of integral ideals of norm $n$
is at most $d(n)$: at a split prime the coefficient at its $a$th
power is $a+1$, while inert and ramified primes give no more.
Multiplicativity proves the bound for every $n$.
Every ideal class here has a representative of norm at most eighteen,
and therefore
\[
              h\bigl(\mathbb Q(\sqrt{-807})\bigr)
                     \leq\sum_{n=1}^{18}d(n)=58<269.
 \tag{PCH.10}
\]
This proves the needed $p$-primality without asserting its exact
class number.

For each nonidentity $h\in\operatorname{Gal}(N/\mathbb Q)=C_2^2$,
norm to its quadratic fixed field followed by extension acts as
$1+h$ on $\operatorname{Cl}(N)[p]$ and is zero.
Thus every nonidentity $h$ acts by minus one.
Two distinct such elements have product the third, which would
then act both by plus one and minus one.
As $p$ is odd, this proves
\[
                           \operatorname{Cl}(N)[p]=0.
 \tag{PCH.11}
\]

We can now prove the full class statement needed for all $j$
simultaneously.  Let $r$ be the reflection in the $S_3$ factor
fixing $L_0$.  Norm to $L_0$ and extension gives $1+r=0$ on
$\operatorname{Cl}(M_0)[p]$ by (PCH.5).
Every conjugate reflection also acts as minus one.
Their products generate the cubic rotations, so $C_3$ acts
trivially on this class module.
On the other hand, norm to $N=M_0^{C_3}$ followed by extension
is $1+c+c^2$ and is zero by (PCH.11).
On the remaining trivial cubic sector it is multiplication by three.
Since $p\ne3$, we obtain
\[
                         \operatorname{Cl}(M_0)[p]=0.
 \tag{PCH.12}
\]
This additional conclusion uses both (PCH.5) and (PCH.11);
it is not inferred from the reflection argument alone.

\proofstep{3. Full-local Kummer descent.}
The degree $[M_0:K]=6$ is prime to $p$.
Restriction/inflation, including each local decomposition or inertia
quotient, identifies the residual strict group with
\[
 \begin{gathered}
 e_{\eta^{-1}}\ker\left(V_{M_0,S}\otimes_{\mathbb F_p}k
                   \longrightarrow\mathcal P_{M_0}\otimes_{\mathbb F_p}k\right),\\
 \mathcal P_{M_0}=\bigoplus_{v\mid p}M_{0,v}^\times/M_{0,v}^{\times p},\\
 V_{M_0,S}=\{x\in M_0^\times:p\mid v(x)\text{ for }v\nmid p\}
                                                /M_0^{\times p}.
 \end{gathered}
 \tag{PCH.13}
\]
Indeed over $M_0$ the coefficient is $k(1)$.
The extra quotient action is the twist $\eta$, so invariants of
Kummer cohomology select the inverse character $\eta^{-1}$.
This remains true when $j=0$ and the cover is not faithful.
It does not require adjoining $\mu_p$ to $M_0$.

The full local power condition also makes the valuations at the
primes above $p$ divisible by $p$.
Thus the kernel in (PCH.13) lies in the all-primes valuation group
$V_{M_0}$, whose exact sequence is
\[
 0\longrightarrow\mathcal O_{M_0}^{\times}/\mathcal O_{M_0}^{\times p}
 \longrightarrow V_{M_0}
 \longrightarrow\operatorname{Cl}(M_0)[p]\longrightarrow0.
 \tag{PCH.14}
\]
By (PCH.12), it remains to prove that localization of the global
units is injective on the whole $b=-1$ part.

\proofstep{4. The unique completion and two exact global units.}
The Eisenstein polynomial $X^6-p$ makes the completion of $L_0$
totally ramified of degree six.
The completion $K_p/\mathbb Q_p$ is unramified quadratic.
Their compositum has local degree twelve, equal to $[M_0:\mathbb Q]$.
Hence $M_0$ has one prime above $p$, with absolute ramification
index six and residue degree two.
The element $z$ is a uniformizer and the powers $1,z,\ldots,z^5$
are an integral basis over $\mathcal O_{K_p}$.
For example, the valuations of $a_nz^n$ for $a_n\in K_p$
are distinct modulo six, so integrality is equivalent to
integrality of all six coefficients.

Two units already contained in the exact AC-G152-class data are
\[
 \begin{aligned}
 U_0&=-82-5z^3,\\
 U_1&=485732-142367z-19217z^2\\
    &\quad+29620z^3-8682z^4-1171z^5.
 \end{aligned}
 \tag{PCH.15}
\]
The first is source unit one, equal to $(279-10w^2)/13$.
The second is source unit four, whose ordered tensor coefficients
in $t^aw^b$, with index $4a+b$, are
\[
 \frac1{3497}\begin{aligned}[t]
 (&-444728016,0,15935560,0,\\
  &17532882,0,-629998,0,\\
  &7935769,0,-284734,0).
 \end{aligned}
 \tag{PCH.16}
\]
Here $3497=13p$.  Substituting the exact formulas (PCH.4) and
$w^2=(p+13z^3)/2$ into the complete rational expression gives
the integral polynomial $U_1$ in (PCH.15).
No term with a denominator divisible by $p$ is discarded.

For a standalone verification, the units satisfy respectively
\[
 \begin{aligned}
 f_0(X)&=X^2+164X-1,\\
 f_1(X)&=X^6-2914392X^5+2292660952446X^4\\
       &\quad+496180748942476X^3-1175165153940X^2\\
       &\quad+13982965350X-1.
 \end{aligned}
 \tag{PCH.17}
\]
Exact Horner evaluation in $\mathbb Q[z]/(z^6-p)$ verifies both
identities.  Their monicity proves integrality, and their constant
coefficient minus one expresses each inverse as an integral
polynomial in the corresponding unit.
The verifier also checks that the source degree-twelve monic
relations are exactly $f_0^6$ and $f_1^2$, respectively.
In addition $U_0b(U_0)=-1$ exactly.
No assertion that these are fundamental units is used.

\proofstep{5. A local witness for each simple unit component.}
Reduce only the integral expressions (PCH.15), obtaining
\[
 \begin{array}{c|rrrrrr}
       &1&z&z^2&z^3&z^4&z^5\\ \hline
 U_0&187&0&0&264&0&0\\
 U_1&187&203&151&30&195&174.
 \end{array}
 \tag{PCH.18}
\]
For each unit, divide by its Teichmuller residue before taking
the local logarithm.  Replacing that residue by the displayed
constant term changes nothing below degree six, since their
difference is divisible by $p=z^6$.
In the ring $\mathbb F_{269^2}[z]/(z^6)$ the resulting logarithms are
\[
 \begin{aligned}
 \ell_0&=128z^3,\\
 \ell_1&=237z+34z^2+222z^3+243z^4+125z^5.
 \end{aligned}
 \tag{PCH.19}
\]
These may be computed without a formal infinite series:
solve $U\ell'=U'$ recursively and integrate, using the invertibility
of the integers one through five modulo $p$.

The first unit is fixed by $c$.
Its actual global quotient $U_0/b(U_0)$ has local logarithm
\[
                         (1-b)\ell_0=256z^3\ne0.
 \tag{PCH.20}
\]
For the second unit first form the actual global cubic correction
$V_1=U_1^3/(U_1c(U_1)c^2(U_1))$.
By (PCH.7), the operator $3-(1+c+c^2)$ kills exponents divisible
by three and multiplies the others by three.
Its logarithm is $173z+102z^2+191z^4+106z^5$.
Applying $1-b$ then gives the logarithm of $V_1/b(V_1)$:
\[
                (1-b)\bigl(3-(1+c+c^2)\bigr)\ell_1
                             =77z+212z^5\ne0.
 \tag{PCH.21}
\]
These are global group-ring operations before localization.
The first class has trivial cubic character and $b=-1$;
the second lies in the nontrivial cubic part with $b=-1$.

The ramification index six is smaller than $p-1$, so local
logarithm and exponential are inverse on the principal units.
A $p$th power principal unit has logarithm of valuation at least
$v_z(p)+1=7$.
The terms of degrees three and one in (PCH.20)--(PCH.21)
therefore give two nonzero full local Kummer classes.
Teichmuller units are themselves $p$th powers, and a root of a
principal unit has residue one because Frobenius is bijective
on the finite residue field.
Thus this test uses the full local power subgroup.

\proofstep{6. Unit multiplicities and all three residual rows.}
The six archimedean places above a fixed complex embedding of $K$
form a regular $\Delta$-set.
Equivariant Dirichlet theory gives the rational unit character
$\mathbb Q[\Delta]-\mathbb Q$.
There are no roots of unity of order $p$ in $M_0$, because
$268$ does not divide $[M_0:\mathbb Q]=12$.
Since $p\nmid6$, reduction modulo $p$ preserves these multiplicities.
The $b=-1$ unit space consequently has dimension three over
$\mathbb F_p$: its cubic-trivial part is one-dimensional, and its
nontrivial cubic part is a simple two-dimensional module.
The latter is simple because $X^2+X+1$ is irreducible over
$\mathbb F_{269}$, as $269\equiv2\pmod3$.

Localization is $\Delta$-equivariant.
The two nonzero values (PCH.20)--(PCH.21) make it injective on
each of these two simple components, and hence on their direct sum.
After extension to $k$, this includes all three inverse-character
lines from (PCH.8).
Together with (PCH.12)--(PCH.14), it proves the residual strict
vanishing for $j=0,1,2$ with the full local conditions.

\proofstep{7. Lift over the original coefficient ring.}
Let $\varpi$ be a uniformizer of the original $O$.
The module $D$ in (PCH.2) is divisible and $D[\varpi]\simeq k(a)$.
For $j=1,2$, its residual cubic inertia at $p$ has order three
and cannot be cancelled by a cyclotomic character of order dividing
$268$.  For $j=0$, the residual inertia character is
$\bar c_p^{-133}$, of exact order $268$, since
$\gcd(133,268)=1$ and $K_p/\mathbb Q_p$ is unramified.
In either case some inertia element acts by a scalar whose
difference from one is a unit of $O$.
It follows that $D^{G_{K_p}}=0$.

Away from $p$ the cyclotomic, Teichmuller and $\kappa$ factors
are unramified.  The remaining inertia character has order
dividing three, prime to $p$.
For every such inertia group the invariants are therefore either
zero or the whole divisible module $D$.
For every local group $J$ in (PCH.1),
$D^J/\varpi D^J=0$.
The multiplication-by-$\varpi$ long exact sequence gives local
injectivity $H^1(J,D[\varpi])\to H^1(J,D)$.

A global strict class killed by $\varpi$ lifts through the global
long exact sequence to a residual class; local injectivity makes
the lift strict.  It is zero by Step6.
Every continuous cocycle from compact $G_K$ to the discrete
module $E/O$ has finite image, annihilated by some power of
$\varpi$.  A nonzero strict class would then give a nonzero one
killed by $\varpi$, a contradiction.
This proves (PCH.2) over the original, possibly ramified ring $O$.
\end{proof}

\paragraph{Actual finite labels and the exact additional weights.}
For $m=0,\kappa=1$, the finite labels are $\chi^{-j}$ for
$j=0,1,2$.  They occur on every actual host containing $K(t)$.
The added residue class consists of $e=133+268n$, $n\geq0$.
For positive integer $T$, its number of pairs $(j,e)$ with
$0<e<T$ is exactly
\[
                         3\left\lfloor\frac{T+134}{268}\right\rfloor.
 \tag{PCH.22}
\]
Together with the previously proved pure cubic rows
$e\equiv0,267\pmod{268}$, the union is
$e\equiv0\pmod{268}$ or $e\equiv-1\pmod{134}$, with count
\[
             3\left(\left\lfloor\frac{T-1}{268}\right\rfloor
                     +\left\lfloor\frac T{134}\right\rfloor\right).
 \tag{PCH.23}
\]
This paragraph records the strict rows and their count only.
The positive derived indices retain their local comparison cost;
the common-window transport requires its separate argument.

\paragraph{Evidence and geometric scope.}
The exact standalone calculation is
\path{verify_pure_cubic_half_units.py}, with output
\path{pure_cubic_half_unit_certificate.json} and an unchanged copy
of the AC-G152-class seven-unit input.
It checks the two rational descents, all original denominators,
both monic unit relations, both projected local values, and the
finite ideal-count bound in (PCH.10).
Its output is a finite unit/local verification, not a new analytic
class certificate.  The full class statement (PCH.12) is proved
above from its precisely stated inherited and elementary inputs.
The source scopes retain Kummer theory from the previously audited
Stacks tag~03PK and the class extension/norm and equivariant
Dirichlet inputs.  SGA~1, Expos\'e~V, \S\S4 and 7 provides the
retained finite-cover/rank-one Artin--Tate dictionary only.
No automorphic endpoint, mixed realization, all-prime estimate,
H04.4 conclusion or Lean execution is asserted.
